\documentclass[11pt]{article}
\usepackage[margin=1in]{geometry}
\usepackage{amsmath,amssymb,amsthm}
\usepackage[hidelinks]{hyperref}

\newtheorem*{theorem}{Literature Answer}

\title{Positive Literature Answer to the \(L_2\) Nonnegative Schauder-Basis Question}
\author{FA/Banach run packet}
\date{2026-07-19}

\begin{document}
\maketitle

\section*{Status}

This is a literature-already-answered packet for a partial subcase.  It records
that the \(p=2\) case of a nonnegative Schauder-basis question mentioned by
Nikolski--Volberg \cite{NV} was answered positively by Freeman--Powell--Taylor
\cite{FPT}.  It is not an original proof from this run.

\section*{Original Problem}

Nikolski--Volberg \cite{NV} construct a normalized uniformly minimal complete
sequence of nonnegative functions in \(L^2_{\mathbb R}(0,1)\), and observe that
this sequence is not a basis.  Immediately afterward they write that the
existence of a nonnegative Schauder basis in \(L^p\), \(p>1\), is open to their
knowledge, while the case \(p=1\) is known by Johnson--Schechtman.

\section*{Supporting Answer}

Freeman--Powell--Taylor \cite{FPT} return directly to this problem.  Their
introduction says that, after the Johnson--Schechtman \(L_1\) construction, the
conditional positive Schauder-basis problem remained open for every
\(1<p<\infty\).  Their main theorem in Section 2 proves:

\begin{theorem}
For every \(\varepsilon>0\), \(L_2(\mathbb R)\) has a Schauder basis consisting
of nonnegative functions, with basis constant at most \(1+\varepsilon\).
\end{theorem}

The same paper remarks that classification theorems extend this conclusion to
all separable \(L_2(\mu)\) spaces.  Therefore the Hilbert-space case of the
Nikolski--Volberg \(p>1\) question has a positive literature answer.

\section*{Scope Limitations}

This packet does not claim a full solution for all \(p>1\).  Freeman--Powell--
Taylor explicitly leave the remaining cases \(1<p<\infty\), \(p\neq 2\), open
as Problem 5.1.  They do prove, for all \(1<p<\infty\), a weaker positive result:
there is a nonnegative Schauder basic sequence whose closed span contains an
isomorphic copy of \(L_p(\mathbb R)\).

\section*{Verification Notes}

The original source and the supporting source are distinct.  The supporting
paper is a later work whose title and Theorem 2.2 address exactly the
nonnegative Schauder-basis problem in the \(L_2\) case.  The local duplicate
search found a failed attempt on separate Markushevich-basis questions from
\cite{FPT}, but no packet recording this partial literature answer to
\cite{NV}.

\begin{thebibliography}{9}

\bibitem{NV}
N. Nikolski and A. Volberg,
\emph{On the Sign Distributions of Hilbert Space Frames},
arXiv:1812.06313.

\bibitem{FPT}
D. Freeman, A. M. Powell, and M. A. Taylor,
\emph{A Schauder basis for \(L_2\) consisting of non-negative functions},
arXiv:2003.09576; Math. Ann. 383 (2022), 181--208.

\end{thebibliography}

\end{document}
